\documentclass[11pt,a4paper]{article}
\usepackage{fontspec,amsmath,amssymb,amsthm,microtype}
\setmainfont{lmroman10-regular.otf}[BoldFont=lmroman10-bold.otf,ItalicFont=lmroman10-italic.otf,BoldItalicFont=lmroman10-bolditalic.otf]
\setmonofont{lmmono10-regular.otf}[Scale=MatchLowercase]
\usepackage[margin=25mm]{geometry}
\usepackage[unicode,hidelinks]{hyperref}
\setlength{\emergencystretch}{3em}
\setlength{\parskip}{0.4em}
\setlength{\parindent}{0pt}
\allowdisplaybreaks[2]
\raggedbottom
\title{\texorpdfstring{The actual frozen random Hessian and its Taylor remainder}{The actual frozen random Hessian and its Taylor remainder}}
\author{LangevinLean mathematical exposition}
\date{}
\begin{document}\maketitle
\noindent Source module: \texttt{Kinetic\allowbreak{}Langevin/\allowbreak{}Lp\allowbreak{}Hessian.\allowbreak{}lean}. The accompanying HTML pairs this mathematical exposition with the exact Lean statements and records the scope of the formalization.


The Hessian in KLMC2 is evaluated at the actual random starting position. It acts on the original coupled L\ensuremath{{}^{2}} variables through the previously constructed random-operator map. Combining its pointwise Taylor remainder with actual fourth displacement moments gives Proposition 2’s quadratic Taylor forcing.


\section{\texorpdfstring{Freeze the true Hessian at a random position}{Freeze the true Hessian at a random position}}


The field \ensuremath{\omega}\ensuremath{\mapsto}H(Y(\ensuremath{\omega})) is strongly measurable because H is Lipschitz and Y has a strongly measurable L\ensuremath{{}^{2}} representative. The force Lipschitz bound controls its derivative: \ensuremath{\Vert}H(y)\ensuremath{\Vert}\ensuremath{\le}M follows from HasFDerivAt g (H(y)) y and the global Lipschitz property of g. The actual random-operator construction then produces a continuous linear action on L\ensuremath{{}^{2}}. Its representative is exactly H(Y(\ensuremath{\omega}))X(\ensuremath{\omega}), and its operator norm is at most M.



{\raggedright\noindent Lean declarations: \texttt{lp\allowbreak{}Hessian}, \texttt{lp\allowbreak{}Hessian\_\allowbreak{}coe}, \texttt{lp\allowbreak{}Hessian\_\allowbreak{}norm\_\allowbreak{}le}.\par}

\section{\texorpdfstring{Lift the genuine Taylor estimate}{Lift the genuine Taylor estimate}}


For actual X,Y, expand the representatives of gLp(X)\ensuremath{-}gLp(Y)\ensuremath{-}H(Y)(X\ensuremath{-}Y). Almost everywhere this is the ordinary pointwise Taylor remainder of g. The deterministic Hessian remainder bound is M\ensuremath{{}_{2}}\ensuremath{\Vert}X\ensuremath{-}Y\ensuremath{\Vert}\ensuremath{{}^{2}}/\allowbreak{}2. If the displacement has fourth norm at most C, the quadratic-to-L\ensuremath{{}^{2}} estimate gives a remainder norm at most M\ensuremath{{}_{2}}C\ensuremath{{}^{2}}/\allowbreak{}2. Every variable stays on the same probability space.



{\raggedright\noindent Lean declarations: \texttt{lp\allowbreak{}Hessian\_\allowbreak{}taylor\_\allowbreak{}bound}.\par}

\section{\texorpdfstring{Obtain the exact dimension factor from stationary velocities}{Obtain the exact dimension factor from stationary velocities}}


The kinematic equation and actual Gaussian velocity marginals give C=t(p\ensuremath{{}^{2}}+2p)\textasciicircum{}(1/\allowbreak{}4). Squaring this radius yields t\ensuremath{{}^{2}}\ensuremath{\surd}(p\ensuremath{{}^{2}}+2p). Consequently the actual L\ensuremath{{}^{2}} Taylor forcing is at most M\ensuremath{{}_{2}}t\ensuremath{{}^{2}}\ensuremath{\surd}(p\ensuremath{{}^{2}}+2p)/\allowbreak{}2. Lean verifies the fractional-power identity using nonnegativity and the real power multiplication rule. No independence of position, Hessian, or velocity is assumed.



{\raggedright\noindent Lean declarations: \texttt{lp\allowbreak{}Hessian\_\allowbreak{}stationary\_\allowbreak{}taylor\_\allowbreak{}bound}.\par}

\end{document}
